\documentclass[11pt]{article}
\usepackage[margin=1.1in]{geometry}
\usepackage{amsmath,amssymb,amsthm}
\usepackage[colorlinks=true,linkcolor=blue,citecolor=blue]{hyperref}

\newtheorem{theorem}{Theorem}[section]
\newtheorem{lemma}[theorem]{Lemma}
\newtheorem{corollary}[theorem]{Corollary}
\newtheorem{proposition}[theorem]{Proposition}
\theoremstyle{definition}
\newtheorem{definition}[theorem]{Definition}
\theoremstyle{remark}
\newtheorem{remark}[theorem]{Remark}
\newtheorem{question}[theorem]{Question}

\newtheorem*{thmA}{Theorem A}
\newtheorem*{thmB}{Theorem B}
\newtheorem*{thmC}{Theorem C}
\newtheorem*{thmD}{Theorem D}
\newtheorem*{thmE}{Theorem E}
\newtheorem*{thmF}{Theorem F}

\newcommand{\Z}{\mathbb{Z}}
\newcommand{\R}{\mathbb{R}}
\newcommand{\ord}{\operatorname{ord}}
\newcommand{\bitlen}{\operatorname{len}}
\newcommand{\LD}{\operatorname{LD}}
\DeclareMathOperator{\vv}{v}

\title{No-Go Theorems for One-Step Lyapunov Potentials\\ for the $3x+1$ Map}
\author{[author]\thanks{Repository: \texttt{github.com/docbgm2002/collatz-things}.
Every numbered claim is machine-verified in exact integer or rational
arithmetic; see Section~\ref{sec:verification}.}}
\date{Draft of November 3, 2026}

\begin{document}
\maketitle

\begin{abstract}
A natural strategy for the $3x+1$ (Collatz) problem is to exhibit a
potential $\Phi(x)=\log_2x+g(x)$ that is nonincreasing along the
accelerated odd map $f(x)=(3x+1)/2^{\vv_2(3x+1)}$, with a correction $g$
computed from simple features of $x$. It is folklore that no such
monotone invariant is known; we prove that for broad natural classes of
corrections none can exist. The master mechanism is \emph{shadowing an
expanding rational cycle}: the cycle $-5\mapsto-7\mapsto-5$ has
multiplier $3^2/2^3>1$, and its positive-integer shadows
$x\equiv-5\pmod{2^N}$ traverse its coordinate itinerary while gaining
value, so any correction that cannot separate $x$ from
$f^2(x)=\tfrac98x+\tfrac58$ is refuted by two integers --- the smallest
certificate being $1275\mapsto1913\mapsto1435$. This excludes all
corrections $g\bigl(x\bmod2^m,\ \tau(x),\ \bitlen(x),\
\lambda_1(x),\dots,\lambda_d(x)\bigr)$, where $\tau$ is the
trailing-ones count and the $\lambda_i$ are suffix-pattern detectors of
arbitrary depth (Theorem~A). Iterating the shadow and applying
Dirichlet approximation to $\log_2\frac98$ extends the exclusion to
corrections using any fixed number of leading binary digits
(Theorem~B); a Mersenne burn argument excludes all bounded corrections
of arbitrary dependence (Theorem~C). The detectors are not ad hoc: they
read a canonical structure, an exact ancestry tower
$w_d(M)=\bigl(2^{M+2d}-2^{2d+1}+3^{d}\bigr)/3^{d}$ with
$f^{d}(w_d(M))=2^{M}-1$, supported on $M\equiv1\pmod{2\cdot3^{d-1}}$
with a carry-free periodic binary normal form (Theorem~D). The
remaining natural coordinate, the size of $x$ read to a fixed
precision $\lfloor2^{j}\log_2x\rfloor$, falls as well: closed
certificates assembled from several shadows of rational cycles exclude
it at every precision $j$, jointly with residues modulo an arbitrary
$N$ and all detectors (Theorem~E), and the least number of witnesses
needed at precision $2^{j}$ is a semiconvergent denominator of
$\log_23$ ($2,2,2,7,7,12,12,12,53,665,\dots$). The exclusions persist
for potentials that are nonincreasing only along $f^{k}$. The same
cycle also blocks a second method: on the Mersenne spine $2^{n}-1\mapsto(3^{n}-1)/2$,
no finite set of residue-class certificates on $n$ can force a positive
block score, because a $2$-adic index $n^{*}$ with $3^{n^{*}+3}=-367$
sends the orbit into $-5\mapsto-7$ (Theorem~F).
\end{abstract}

\section{Introduction}

For odd $x\in\Z_{>0}$ let
\[
f(x)=\frac{3x+1}{2^{\vv_2(3x+1)}},
\]
the accelerated (odd-to-odd) form of the Collatz map; the $3x+1$
conjecture asserts that every orbit reaches $1$. Among the strategies
proposed for the problem, one of the most natural is the
\emph{potential-function} (Lyapunov) method: find
$\Phi:2\Z+1\to\R$ with $\Phi(f(x))\le\Phi(x)$ for all odd $x>1$,
ideally of the coercive form $\Phi(x)=\log_2x+g(x)$ with $g$ a
correction of lower order, so that monotonicity constrains orbit
growth. Heuristics support the idea: a ``random'' odd step multiplies
$x$ by $3/2^{e}$ with $e$ geometrically distributed of mean $2$, so
$\log_2x$ decreases by $\log_2\frac43$ \emph{on
average}~\cite{Lagarias1985}. The difficulty, as surveys note, is that
no correction making the decrease pointwise has ever been
found~\cite{Lagarias2010}. To the best of our knowledge the literature
contains no theorem \emph{excluding} a concrete class of such
potentials; the absence is reported as an empirical fact. This paper
proves exclusions, and locates the precise boundary at which the
exclusion program must stop.

Throughout, $\tau(x)=\vv_2(x+1)$ is the number of trailing
\texttt{1}s in the binary expansion of odd $x$ (the \emph{fuel}:
Lemma~\ref{lem:fuse} shows an odd step converts one unit of fuel into a
multiplication by almost exactly $\tfrac32$); $\bitlen(x)$ is the bit
length; $\LD_j(x)=\lfloor x/2^{\bitlen(x)-j}\rfloor$ is the top $j$
bits; and $\lambda_1,\lambda_2,\dots$ are the suffix-pattern detectors
of Definition~\ref{def:detector}, which read the run length of a
canonical hierarchy of periodic binary patterns (levels $1$ and $2$
being $(\texttt{10})^{k}\texttt{1001}$ and
$(\texttt{111000})^{k}\texttt{01}$). Our main impossibility results:

\begin{thmA}[master no-go; Theorem~\ref{thm:A}]
For every $m,d\ge0$ and every function $g$, the potential
$\log_2x+g\bigl(x\bmod2^m,\ \tau(x),\ \bitlen(x),\
\lambda_1(x),\dots,\lambda_d(x)\bigr)$ is not nonincreasing along $f$
on odd $x>1$.
\end{thmA}

\begin{thmB}[leading digits; Theorem~\ref{thm:B}]
For every $m,j,d\ge0$ and every $g$, the potential
$\log_2x+g\bigl(x\bmod2^m,\ \tau(x),\ \LD_j(x),\
\lambda_1(x),\dots,\lambda_d(x)\bigr)$ is not nonincreasing along $f$
on odd $x>1$.
\end{thmB}

\begin{thmC}[bounded corrections; Theorem~\ref{thm:C}]
For every bounded $G:2\Z+1\to\R$, of arbitrary dependence on $x$, the
potential $\log_2x+G(x)$ is not nonincreasing along $f$ on odd $x>1$.
\end{thmC}

\begin{thmE}[quantized logarithm; Theorem~\ref{thm:E}]
For every $j,m,d\ge0$, every $N\ge1$ and every $g$, the potential
$\log_2x+g\bigl(\lfloor2^{j}\log_2x\rfloor,\ \tau(x),\ x\bmod N,\
\lambda_1(x),\dots,\lambda_d(x)\bigr)$ is not nonincreasing along $f$
on odd $x>1$; nor is it nonincreasing along $f^{k}$ for any fixed $k$.
\end{thmE}

Theorem~A is proved by two integers per parameter choice: the map has
the expanding rational cycle $-5\mapsto-7\mapsto-5$, of multiplier
$3^2/2^3=\tfrac98>1$ (expansion is exactly why it lives at negative
integers), and its shadow at positive $x\equiv-5\pmod{2^N}$ reproduces
the cycle's coordinate itinerary while the value gains a factor
$>\tfrac98$ per period. Any correction blind to the difference between
$x$ and $f^2(x)=\tfrac98x+\tfrac58$ is refuted outright; the smallest
certificate is $1275\mapsto1913\mapsto1435$. Theorem~B iterates the
shadow $a$ times and chooses $a$, by Dirichlet approximation to the
irrational $\log_2\tfrac98$, so that the accumulated gain sits next to
an exact power of $2$ and the leading digits return. Theorem~C is a
one-paragraph consequence of the Mersenne burn. Theorem~E replaces the
single shadow by a closed cycle of $K$ independent witnesses, each
shadowing one member of a rational cycle chosen to fit the precision,
and the number $K$ needed is exactly the least denominator at which
$\lfloor K\log_23\rfloor/K$ under-approximates $\log_23$ at least as
well as $\lfloor2^{j}\log_23\rfloor/2^{j}$ does (Section~\ref{sec:quantized}).

The detectors quantified over in Theorems A and B are given meaning by
a structural result of independent interest:

\begin{thmD}[Mersenne ancestry tower; Theorem~\ref{thm:D}]
For $d\ge1$ and $M\equiv1\pmod{2\cdot3^{d-1}}$, $M>1$, the integer
$w_d(M)=\bigl(2^{M+2d}-2^{2d+1}+3^{d}\bigr)/3^{d}$ is odd with
$\tau=1$, satisfies $3w_d(M)+1=4\,w_{d-1}(M)$ and hence
$f^{d}(w_d(M))=2^{M}-1$ with every step of value ratio in
$(\tfrac34,1)$; its binary expansion is exactly periodic (carry-free)
of period $2\cdot3^{d-1}$; and $w_d(M)\equiv1-2^{2d+1}3^{-d}
\pmod{2^{M-1}}$, distinct levels separating $2$-adically at bit
$2\min(d,d')+1$ exactly.
\end{thmD}

The tower is the complete uniform-$e{=}2$ ancestry of the Mersenne
numbers --- the slowest exact approach to maximal fuel --- and
$\lambda_d$ detects membership in its level-$d$ suffix pattern. The
congruence domain is the order of $2$ modulo $3^d$, and the carry-free
normal form is an instance of the lifting-the-exponent lemma.

Section~\ref{sec:spine} turns from potentials to congruence case
analysis. After the burn, $2^{n}-1$ reaches the repunit $a_n=(3^{n}-1)/2$,
and a natural attack on its stopping time scores the orbit of $a_n$ block
by block and splits on residues of $n$. Each such case is an exact finite
computation (Proposition~\ref{prop:decide}), but the method cannot cover
all $n$:

\begin{thmF}[Theorem~\ref{thm:F}]
On $n=64k+17$ with $k$ odd, the $2$-adic index $k^{*}$ defined by
$3^{64k^{*}+20}=-367$ sends the orbit of $a_n$ into the cycle
$-7\mapsto-5\mapsto-7$, every block of which has negative score. Integers
$k\equiv k^{*}\pmod{2^{J}}$ follow it for $\lfloor(J+2)/3\rfloor$ blocks,
and no finite set of certified residue classes covers all odd $k$.
\end{thmF}

\subsection*{Method and provenance}

The constraint ``$\Phi$ nonincreasing'' restricted to any coordinate
window is a difference-constraint system on the values of $g$; it is
infeasible if and only if the coordinate graph of real orbit steps
contains a directed cycle whose product of value ratios exceeds $1$
(Section~\ref{sec:certificates}). Mining such cycles by Bellman--Ford
over odd $x\le10^6$ produces startlingly short certificates --- e.g.\
the pair $57\mapsto43$, $27\mapsto41$ already refutes all
$g(x\bmod16,\tau)$ --- and inspection of the mined witnesses revealed
that they shadow the $-5$ cycle, which led to the uniform theorems
above. We record this provenance because the historical route of this
project ran in the opposite direction: Theorems A--C were first proved
for subclasses by designed families (a Mersenne ``burn'' that forces
linear growth of $g$ along a frozen coordinate line, against a
``recharge'' family that caps it), a method that yields an exact price
list --- each level of suffix detection buys a would-be potential
exactly $\log_2\frac43$ of budget against an unbounded liability ---
but is dominated, as a source of impossibility statements, by the
two-integer shadow certificates. The family analysis survives in the
repository accompanying this paper and in the structural
Theorem~D.

\subsection*{The boundary}

A coordinate defeats a single shadow only by separating $x$ from
$f^{2a}(x)$ for every good Dirichlet approximation $a$; the quantized
logarithm $\lfloor2^{j}\log_2x\rfloor$ does so for $j\ge3$, since one
shadow period gains $\log_2\tfrac98>2^{-j}$. Theorem~E shows that
this refuge is illusory: certificates built from several witnesses
reset their position in the cell at every junction, and exist at every
precision. After Theorems A, B and E, no correction built from finitely
much arithmetic information about $x$ --- its residue modulo any fixed
integer, any fixed-depth suffix pattern, its size to any fixed precision
--- makes $\log_2x$ monotone. Section~\ref{sec:boundary} records what
is left: the leading digits taken jointly with the quantized logarithm,
and corrections that depend on the orbit rather than on $x$, for which
the existence of a monotone potential is conjecture-adjacent.

\subsection*{Relation to the literature}

The parity-vector structure underlying the heuristic goes back to
Terras~\cite{Terras1976} and Everett~\cite{Everett1977}; density
results culminate in~\cite{KrasikovLagarias2003} and, by different
methods, in Tao's almost-all theorem~\cite{Tao2022}. Cycle exclusions
rest on transcendence bounds~\cite{Steiner1977,SimonsDeWeger2005} and
computation~\cite{Barina2021}. Backward (predecessor) structure of the
Collatz graph is developed at the level of density and measure by
Wirsching~\cite{Wirsching1998}; the tower of Theorem~D is a
distinguished explicit chain in that graph, and its closed form,
congruence domain, and $2$-adic freezing appear to be new. Forward
identities of burn type, attached to the rational cycles at $-1$, $-5$,
$-17$, appear in Andaloro~\cite{Andaloro2000} --- the shadows of
Section~\ref{sec:shadow} are those identities viewed through the
certificate lens. A search of the literature, including the annotated
bibliographies~\cite{LagariasBibI,LagariasBibII}, located no prior
impossibility theorems for concrete classes of one-step potentials; we
would welcome references.

\section{Preliminaries}\label{sec:prelim}

Odd integers are written $x$; $\vv_p$ is the $p$-adic valuation;
$\tau(x)=\vv_2(x+1)$; congruential inverses are taken in
$(\Z/2^m\Z)^{\times}$. Two standard facts are used.

\begin{lemma}[LTE, special case]\label{lem:lte}
For even $n$, $\vv_3(2^{n}-1)=1+\vv_3(n)$; in particular
$\vv_3\bigl(2^{2\cdot3^{d-1}}-1\bigr)=d$.
\end{lemma}

\begin{lemma}\label{lem:order}
$2$ is a primitive root modulo $3^d$ for every $d\ge1$; hence
$\ord_{3^d}(2)=\varphi(3^d)=2\cdot3^{d-1}$.
\end{lemma}

\begin{proof}
$2$ generates $(\Z/9\Z)^{\times}$ (its order divides $6$ and is not
$1,2,3$), and a primitive root mod $9$ is a primitive root mod $3^d$
for all $d$.
\end{proof}

\begin{lemma}[fuse]\label{lem:fuse}
If $\tau(x)=t\ge2$ then $\vv_2(3x+1)=1$, $f(x)=\tfrac{3x+1}{2}$,
$\tau(f(x))=t-1$, and $f(x)/x=\tfrac32+\tfrac1{2x}$.
\end{lemma}

\begin{proof}
Write $x+1=2^{t}u$, $u$ odd. Then
$3x+1=2\,(3\cdot2^{t-1}u-1)$ with odd bracket since $t\ge2$, and
$f(x)+1=3\cdot2^{t-1}u$.
\end{proof}

\begin{remark}\label{rem:andaloro}
The iterated form of Lemma~\ref{lem:fuse} --- that $2^{t}u-1$ ($u$ odd)
iterates in $t$ odd steps to $3^{t}u-1$, the \emph{burn} --- is
classical; it is recorded by Andaloro~\cite{Andaloro2000} (cf.\ the
annotation in~\cite{LagariasBibII}).
\end{remark}

Finally we define the detectors. For $d\ge1$ write
$\ord_d=2\cdot3^{d-1}$ and, for $M\equiv1\pmod{\ord_d}$, $M>1$,
\begin{equation}\label{eq:wd}
w_d(M)=\frac{2^{M+2d}-2^{2d+1}+3^{d}}{3^{d}}
=1+2^{2d+1}\,\frac{2^{M-1}-1}{3^{d}} ,
\end{equation}
an odd integer by Lemma~\ref{lem:order} (Theorem~\ref{thm:D} develops
its structure; nothing beyond the formula and Lemma~\ref{lem:mod32}
below is needed before Section~\ref{sec:tower}).

\begin{definition}\label{def:detector}
For $d\ge1$ let $\beta_d(s)$ be the bit length of
$T^{(d)}_s:=w_d(1+\ord_d\,s)$ and define the \emph{level-$d$ detector}
\[
\lambda_d(x)=\max\bigl\{\,s\ge1:\
x\equiv T^{(d)}_s\ \bigl(\mathrm{mod}\ 2^{\,\beta_d(s)+1}\bigr)\bigr\}
\ \in\ \Z_{\ge1}\cup\{\bot\},
\]
the maximum of the empty set being $\bot$.
\end{definition}

\begin{lemma}\label{lem:mod32}
Every template satisfies $T^{(d)}_s\equiv9\pmod{32}$ if $d=1$ and
$T^{(d)}_s\equiv1\pmod{32}$ if $d\ge2$. Consequently
$\lambda_d(y)=\bot$ for every $d$ whenever
$y\not\equiv1,9\pmod{32}$.
\end{lemma}

\begin{proof}
For $d\ge2$, $2^{2d+1}\equiv0\pmod{32}$ in \eqref{eq:wd}. For $d=1$:
$M$ odd, $M\ge3$, so $2^{M-1}\equiv0\pmod4$, whence
$(2^{M-1}-1)/3\equiv(-1)\cdot3^{-1}\equiv1\pmod4$ and
$w_1(M)\equiv1+8\equiv9\pmod{32}$. Matching any template forces the
mod-$32$ congruence, since every matching modulus
$2^{\beta_d(s)+1}$ contains $32$ ($\beta_d(s)\ge4$).
\end{proof}

\section{The master theorem: shadowing the $-5$ cycle}\label{sec:shadow}

The map has the rational cycle $-5\mapsto-7\mapsto-5$:
$f(-5)=-7$, $f(-7)=-5$, with multiplier $3^2/2^3=\tfrac98>1$. Expanding
cycles live at negative integers (the cycle equation gives
$x=c/(2^{E}-3^{K})$ with $2^{E}<3^{K}$), and their positive shadows
gain value while walking the cycle's coordinate itinerary.

\begin{lemma}[shadow period]\label{lem:S}
Let $N\ge8$, $w$ odd, and $x=2^{N}w-5$. Then
\[
\vv_2(3x+1)=1,\quad f(x)=3\cdot2^{N-1}w-7,\qquad
\vv_2(3f(x)+1)=2,\quad f^{2}(x)=9\cdot2^{N-3}w-5,
\]
so that $8\,f^{2}(x)=9x+5$ and $f^{2}(x)/x=\tfrac98+\tfrac5{8x}$.
Moreover $\tau(x)=\tau(f^2(x))=2$ and $\tau(f(x))=1$;
$x\equiv f^{2}(x)\equiv-5$ and $f(x)\equiv-7\pmod{2^{\,N-3}}$; and
$\lambda_d(x)=\lambda_d(f(x))=\lambda_d(f^2(x))=\bot$ for every $d$.
\end{lemma}

\begin{proof}
Direct computation: $3x+1=2(3\cdot2^{N-1}w-7)$ with odd bracket;
$3f(x)+1=4(9\cdot2^{N-3}w-5)$ with odd bracket; $8f^2(x)=9\cdot2^Nw-40
=9x+5$. Fuel: $x+1=4(2^{N-2}w-1)$;
$f(x)+1=2\cdot3(2^{N-2}w-1)$ with odd second factor;
$f^2(x)+1=4(9\cdot2^{N-5}w-1)$ (odd bracket, $N\ge6$). Residues are
immediate. Detectors: $x\equiv f^2(x)\equiv27$ and
$f(x)\equiv25\pmod{32}$ (using $N-3\ge5$), and
$27,25\notin\{1,9\}$; apply Lemma~\ref{lem:mod32}.
\end{proof}

\begin{theorem}[master no-go]\label{thm:A}
Let $m,d\ge0$ and let $g$ be any real-valued function of
$\bigl(x\bmod2^m,\ \tau(x),\ \bitlen(x),\
\lambda_1(x),\dots,\lambda_d(x)\bigr)$. Then
$\Phi(x)=\log_2x+g(\cdots)$ is not nonincreasing along $f$ on odd
$x>1$.
\end{theorem}

\begin{proof}
Take $N\ge\max(m+3,8)$, $k\ge3$, and $w=2^{k-1}+1$ in
Lemma~\ref{lem:S}. Then additionally
$\bitlen(x)=\bitlen(f^{2}(x))=N+k$, since $w$ places $x$ in the lower
part of its dyadic window: $9(2^{k}+2)<2^{k+4}$. Hence
$x$ and $f^{2}(x)$ have identical coordinate vectors while
$f^{2}(x)>x$, and summing the two monotonicity constraints at $x$ and
$f(x)$ gives
$0\le-\log_2\bigl(f^{2}(x)/x\bigr)<-\log_2\tfrac98<0$.
\end{proof}

The smallest member ($N=8$, $k=3$) is
$1275\mapsto1913\mapsto1435$, with $8\cdot1435=9\cdot1275+5$ and ratio
product $\tfrac{287}{255}$: two integers refute every $g$ on the full
coordinate algebra at $m\le5$.

\begin{remark}[general principle]\label{rem:factories}
Nothing is special about $-5$: every expanding rational cycle of the
map --- every solution of the cycle equation with $2^{E}<3^{K}$, e.g.\
the $K=7$ cycle at $-17$ with multiplier $3^{7}/2^{11}$ --- shadows at
positive integers $\equiv$ (cycle member) $\bmod\ 2^{N}$ into a closed
coordinate itinerary of value gain $3^{K}/2^{E}>1$, refuting every
correction that is $2$-adically local along it. Expanding negative
cycles are certificate factories against local potentials.
\end{remark}

\section{Iterating the shadow: leading digits}\label{sec:leading}

A coordinate can defeat a single shadow period only by separating $x$
from $f^{2}(x)=\tfrac98x+\tfrac58$; the leading digits do (the ratio
$\tfrac98$ shifts them). Iteration removes this refuge.

\begin{theorem}[leading digits do not survive]\label{thm:B}
For every $m,j,d\ge0$ and every $g$, the potential
$\log_2x+g\bigl(x\bmod2^m,\ \tau(x),\ \LD_j(x),\
\lambda_1(x),\dots,\lambda_d(x)\bigr)$ is not nonincreasing along $f$
on odd $x>1$.
\end{theorem}

\begin{proof}
For $x\equiv-5\pmod{2^{\,3a+c}}$, $c\ge5$, Lemma~\ref{lem:S} applies
$a$ times (each period consumes three bits of shadow depth), the
itinerary alternating residues $-5,-7$ with fuel $2,1$ and every
detector $\bot$ throughout, giving the exact closed form
\[
f^{2a}(x)=\frac{9^{a}x+5\,(9^{a}-8^{a})}{8^{a}}
=\Bigl(\tfrac98\Bigr)^{a}x\,(1+\delta),\qquad 0<\delta<\tfrac5x .
\]
Now $\beta=\log_2\tfrac98$ is irrational ($2^{\,p+3q}=3^{\,2q}$ is
impossible: $\vv_3$ of the sides differ), so by Dirichlet there is
$a\ge1$ with $\|a\beta\|<2^{-j-3}$ (distance to the nearest integer
$P$; note $a\beta>0$ regardless of the side). Choose
$N=3a+\max(m,5)+j+8$ and a free odd multiplier $w$ placing the
fractional part of $\log_2x$, $x=2^{N}w-5$, at distance $>2^{-j-2}$
from every breakpoint of $\LD_j$ (possible since
$\log_2(2^{N}w-5)\approx N+\log_2w$ and $\{\log_2w\}$ is dense), with
$N$ large enough that $\log_2(1+\delta)<2^{-j-3}$. Then
$\log_2f^{2a}(x)=\log_2x+P\pm\|a\beta\|+\log_2(1+\delta)$, so
$\LD_j$ closes; all other coordinates close by the itinerary; and
telescoping monotonicity over the $2a$-step walk gives
$0\ge a\beta+\log_2(1+\delta)>0$.
\end{proof}

A concrete instance ($j=4$): $a=53$ gives $53\beta=9.0060\ldots$; the
verifier realizes the full $106$-step walk in exact integers at depth
$N=177$, with top-$4$ bits $\texttt{1011}\to\texttt{1011}$, all
coordinates frozen, and value gain $2^{9.006}$.

\section{Bounded corrections}\label{sec:bounded}

\begin{theorem}[no bounded correction]\label{thm:C}
Let $G:2\Z+1\to\R$ be any bounded function. Then $\log_2x+G(x)$ is not
nonincreasing along $f$ on odd $x>1$.
\end{theorem}

\begin{proof}
Let $\Omega=\sup G-\inf G$. Telescoping monotonicity over the burn of
$2^{M}-1$ (Lemma~\ref{lem:fuse} iterated: $M-1$ steps, each of ratio
$>\tfrac32$) forces $(M-1)\log_2\tfrac32\le\Omega$, false for
$M>1+\Omega/\log_2\tfrac32$.
\end{proof}

Theorems A--C partition the correction landscape: bounded corrections
fail on the burn alone; unbounded corrections through residues, fuel,
length, leading digits, and suffix detectors of every depth fail on
the shadow. Section~\ref{sec:quantized} treats the size of $x$ read to
an arbitrary fixed precision, and Section~\ref{sec:boundary} describes
what is left.

\section{The Mersenne ancestry tower}\label{sec:tower}

We now develop the structure behind the detectors. Define
$w_0(M)=2^{M}-1$ and $w_d(M)$ as in \eqref{eq:wd}.

\begin{theorem}\label{thm:D}
For $d\ge1$ and $M\equiv1\pmod{\ord_d}$, $M>1$:
\begin{enumerate}
\item[(a)] \textup{(Domain.)} $3^{d}$ divides
$2^{M+2d}-2^{2d+1}+3^{d}$ iff $M\equiv1\pmod{\ord_d}$.
\item[(b)] \textup{(Form.)} $w_d(M)$ is an odd integer $>1$ with
$\tau(w_d(M))=1$.
\item[(c)] \textup{(Descent.)} $3\,w_d(M)+1=4\,w_{d-1}(M)$; hence
$\vv_2(3w_d(M)+1)=2$, $f(w_d(M))=w_{d-1}(M)$, and
$f^{d}(w_d(M))=2^{M}-1$, each step of value ratio
$\tfrac34+\tfrac1{4w_d}\in(\tfrac34,1)$.
\item[(d)] \textup{(Normal form.)} $B_d:=(2^{\ord_d}-1)/3^{d}$ is an
integer in $(0,2^{\ord_d})$ by Lemma~\ref{lem:lte}, and for
$M=1+\ord_d\,s$,
\[
w_d(M)=1+2^{2d+1}\cdot\bigl[\text{the $\ord_d$-bit block of
}B_d\bigr]^{s},
\]
the bracket denoting $s$ carry-free repetitions: an exactly periodic
binary word of period $2\cdot3^{d-1}$.
\item[(e)] \textup{(Freeze and separation.)}
$w_d(M)\equiv\Lambda_d:=1-2^{2d+1}3^{-d}\pmod{2^{\,M-1}}$; and for
$d\ne d'$, $\vv_2(\Lambda_d-\Lambda_{d'})=2\min(d,d')+1$.
\end{enumerate}
\end{theorem}

\begin{proof}
(a) Modulo $3^{d}$ the numerator is $2^{2d}(2^{M}-2)$, so divisibility
is $2^{M-1}\equiv1\pmod{3^d}$; apply Lemma~\ref{lem:order}.
(b) The second form in \eqref{eq:wd} is $1$ plus an even number, and
$w_d+1=2(1+2^{2d}q)$ with $2^{2d}q$ even.
(c) Clearing $3^{d-1}$, the identity reduces to
$3^{d}+3^{d-1}=4\cdot3^{d-1}$; $w_{d-1}$ is odd, so $\vv_2=2$ exactly;
iterate.
(d) With $M-1=\ord_d\,s$,
$(2^{M-1}-1)/3^{d}=B_d\sum_{i=0}^{s-1}2^{\ord_d i}$, a
base-$2^{\ord_d}$ number all of whose digits equal $B_d<2^{\ord_d}$: no
carries; the $+1$ and the shift by $2^{2d+1}$ cannot interact.
(e) $3^{d}\cdot\tfrac{2^{M-1}-1}{3^{d}}\equiv-1\pmod{2^{M-1}}$; and for
$d>d'$,
$\Lambda_d-\Lambda_{d'}=2^{2d'+1}3^{-d}\bigl(3^{\,d-d'}-2^{\,2(d-d')}\bigr)$
with odd bracket.
\end{proof}

Levels $0,1,2$ are the Mersennes, the family
$x'_M=(2^{M+2}-5)/3=(\texttt{10})^{(M-3)/2}\texttt{1001}$ with
$f(x'_M)=2^M-1$, and $(\texttt{111000})^{k}\texttt{01}$; the tower is
the complete uniform-$e{=}2$ ancestry of maximal fuel, and every tower
member lies on the residue class $8y+1$. Composing the tower with the
burn (Remark~\ref{rem:andaloro}) yields, for each $(d,M)$, an exactly
known orbit segment of length $d+M+1$ terminating on the base-$3$
repunit $(3^{M}-1)/2$.

\begin{lemma}\label{lem:detfacts}
(1) $\lambda_d(T^{(d)}_s)=s$: the detector reads its own tower index,
unboundedly. (2) For $d\ne d'$, with $\nu=2\min(d,d')+1$: for all
$M>\nu+1$ in the level-$d$ progression,
$\lambda_{d'}(w_d(M))\le\nu/\ord_{d'}$.
\end{lemma}

\begin{proof}
(1) At index $s$ the congruence is an equality that holds; at $s'>s$
the modulus exceeds both operands and templates strictly increase.
(2) Suppose $w_d(M)$ matches $T^{(d')}_s$ with $M'=1+\ord_{d'}s$ and,
toward a contradiction, $M'-1\ge\nu+1$. The matching modulus refines
$2^{\nu+1}$ (as $\beta_{d'}(s)\ge M'$), so Theorem~\ref{thm:D}(e) on
both sides forces $\Lambda_d\equiv\Lambda_{d'}\pmod{2^{\nu+1}}$,
contradicting $\vv_2(\Lambda_d-\Lambda_{d'})=\nu$. Hence
$\ord_{d'}s\le\nu$. (Small-index templates have not converged to
$\Lambda_{d'}$; this affects which small values occur, not the bound.)
\end{proof}

\begin{remark}[the family method and the price list]\label{rem:price}
Historically, Theorems A--C were first obtained for subclasses by a
family method that yields quantitative structure the shadow does not.
The \emph{burn} family $3^{j}2^{t}-1$ ascends at ratio $>\tfrac32$ per
step with frozen residue $-1$ and $\lambda\equiv\bot$, forcing $g$ to
grow at slope $\log_2\tfrac32$ along a coordinate line; the tower's
$e{=}2$ chains cap the same values from above, each level costing one
$\log_2\tfrac43$. The resulting price list --- each detector level buys
a would-be potential exactly $\log_2\tfrac43$ of budget against an
unbounded liability --- explains \emph{why} suffix information fails,
while the shadow certificates prove \emph{that} it fails. The full
family proofs are retained in the repository.
\end{remark}

\section{The quantized logarithm}\label{sec:quantized}

Theorems A and B leave one natural coordinate untouched: the size of
$x$ read to a fixed precision, integer and fractional parts of
$\log_2x$ jointly. Write $Q=2^{j}$ and
\[
c(x)=\lfloor Q\log_2x\rfloor ,
\]
the \emph{quantized logarithm} at precision $j$ ($c(x)=\bitlen(x)-1$
for $j=0$). A single shadow chain cannot close it: one period of the
shadow gains $\log_2\tfrac98>2^{-j}$ for $j\ge3$, so $x$ and
$f^{2a}(x)$ never share a cell. The remedy is to use several witnesses,
each a fresh integer: at every junction the position inside the cell is
reset, and the total displacement around the cycle can be made zero
exactly. This section proves that such certificates exist at every
precision, for arbitrary moduli and detectors at once, and that their
minimal length is a Diophantine invariant of $\log_23$.

Throughout, $\beta=\log_23$, $L=\lfloor Q\beta\rfloor$,
$\varphi=Q\beta-L\in(0,1)$ (irrational $\beta$), and for odd $x$
\[
\theta(x)=Q\log_2x-c(x)\in[0,1),\qquad
\eta(x)=Q\log_2\Bigl(1+\frac1{3x}\Bigr)\in\Bigl(0,\frac{Q}{3x\ln2}\Bigr).
\]

\begin{definition}\label{def:certificate}
Let $C$ be any function on the odd integers (a \emph{coordinate}). A
\emph{closed certificate} of length $K$ for $C$ is a list of odd
integers $x_1,\dots,x_K>1$, not necessarily distinct, with
\[
C(f(x_i))=C(x_{i+1})\quad(i\bmod K)\qquad\text{and}\qquad
\prod_{i=1}^{K}f(x_i)>\prod_{i=1}^{K}x_i .
\]
\end{definition}

\begin{lemma}\label{lem:certrefutes}
If a closed certificate exists for $C$, then no potential
$\log_2x+g(C(x))$ is nonincreasing along $f$ on odd $x>1$.
\end{lemma}

\begin{proof}
Sum the $K$ inequalities $\Phi(f(x_i))\le\Phi(x_i)$. The $g$-terms
cancel in pairs by closure, leaving
$\sum_i\log_2\bigl(f(x_i)/x_i\bigr)\le0$, which contradicts the gain.
\end{proof}

The number of witnesses is governed by the following quantity.

\begin{definition}\label{def:admissible}
$K\ge1$ is \emph{admissible at precision $Q$} if
\[
s(Q,K):=Q\lfloor K\beta\rfloor-KL\ \ge\ 0,
\qquad\text{equivalently}\qquad
\frac{\lfloor K\beta\rfloor}{K}\ \ge\ \frac{\lfloor Q\beta\rfloor}{Q}.
\]
\end{definition}

The identity $s(Q,K)=K\varphi-Q\{K\beta\}$ shows $s(Q,K)<K$ always,
and $s(Q,Q)=0$, so $K=Q$ is admissible at every precision.

\subsection*{Two lemmas}

\begin{lemma}[displacement]\label{lem:displacement}
Let $x$ be odd with $e=\vv_2(3x+1)$. Then
\[
c(f(x))=c(x)+L-Qe+\delta,\qquad
\delta=\bigl\lfloor\theta(x)+\varphi+\eta(x)\bigr\rfloor .
\]
In particular, if $\eta(x)\le1-\varphi$ then $\delta\in\{0,1\}$, with
$\delta=0$ whenever $\theta(x)<1-\varphi-\eta(x)$ and $\delta=1$
whenever $\theta(x)\ge1-\varphi$. The last condition is the exact
integer inequality $(3x)^{Q}\ge2^{\,c(x)+1+L}$.
\end{lemma}

\begin{proof}
$Q\log_2f(x)=Q\log_2x+Q\beta-Qe+\eta(x)
=c(x)+L-Qe+\bigl(\theta(x)+\varphi+\eta(x)\bigr)$.
For the integer form: $(3x)^{Q}\ge2^{c+1+L}$ iff
$Q\beta+c+\theta\ge c+1+L$ iff $\theta\ge1-\varphi$.
\end{proof}

\begin{lemma}[rational cycles]\label{lem:ratcycle}
Let $(e_1,\dots,e_K)$ be any word of positive integers with sum $E$,
and put $E_i=e_1+\dots+e_i$. Then
\[
y_1=\frac{\sum_{i=1}^{K}3^{K-i}\,2^{E_{i-1}}}{2^{E}-3^{K}}
\]
is a rational number whose denominator is prime to $6$; its orbit under
$f$ (extended to $\Z_2$) is purely periodic, $y_{i+1}=(3y_i+1)/2^{e_i}$
and $y_{K+1}=y_1$, with $\vv_2(3y_i+1)=e_i$ exactly for every $i$. Every
$y_i$ is a $2$-adic unit and a $3$-adic integer.
\end{lemma}

\begin{proof}
Composing the affine maps $y\mapsto(3y+1)/2^{e_i}$ gives
$2^{E}y_{K+1}=3^{K}y_1+\sum_i3^{K-i}2^{E_{i-1}}$, so $y_1$ as stated is
the unique fixed point of the composite; the same formula applied to the
rotated words shows that every $y_i$ has the shape
$c_i/(2^{E}-3^{K})$ with $c_i\in\Z$, so every $y_i\in\Z_2\cap\Z_3$. If
some $y_{i+1}$ were even, then $3y_{i+1}+1$ would be a $2$-adic unit and
$y_{i+2}=(3y_{i+1}+1)/2^{e_{i+1}}\notin\Z_2$, a contradiction. Hence
every $y_{i+1}$ is odd, i.e.\ $\vv_2(3y_i+1)=e_i$. (Periodic parity
sequences of rational $2$-adic integers are classical; see
Lagarias~\cite{Lagarias1985}, Theorem~L.)
\end{proof}

For the word $(1,2)$ the cycle is $-5\mapsto-7$; for $(2)$ it is the
fixed point $1$. Lemma~\ref{lem:S} is the shadow of $(1,2)$; the proof
below shadows the cycle of a word chosen to fit the precision.

\subsection*{The theorem}

\begin{theorem}[the quantized logarithm does not survive]\label{thm:E}
Let $j,m,d\ge0$ and $N\ge1$, and let $g$ be any real-valued function
of
\[
\bigl(\lfloor2^{j}\log_2x\rfloor,\ \tau(x),\ x\bmod N,\
\lambda_1(x),\dots,\lambda_d(x)\bigr).
\]
Then $\log_2x+g(\cdots)$ is not nonincreasing along $f$ on odd $x>1$.
More precisely, write $N=2^{m}3^{u}N'$ with $\gcd(N',6)=1$. For every
$K\ge2$ admissible at $Q=2^{j}$ with $2\lfloor K\beta\rfloor\ge3K$ ---
a condition satisfied by every admissible $K\ge2$ when $j\ge1$ --- there
is a closed certificate for the coordinate of length $tK$, where $t\ge1$
depends only on $K$ and $N'$ and $t=1$ when $N'=1$. Its witnesses can be
taken of bit length at most
\[
\Theta=m+\log_2(3^{u}N')+6tK+j+\log_2\frac1{\min(\varphi,1-\varphi)}+10 .
\]
\end{theorem}

\begin{proof}
Put $E=\lfloor K\beta\rfloor$ and $s=s(Q,K)\in[0,K)$.

\emph{Step 1: the word.} Let $p=5K-3E$ and $r=2E-3K$. Both are
nonnegative ($\beta<\tfrac53$ and the hypothesis; for $j\ge1$ the
hypothesis follows from admissibility, since $L/Q\ge\tfrac32$ for
$Q=2,4,8$ and $Q\beta-1\ge\tfrac32Q$ for $Q\ge12$), and the cyclic word
\[
(e_1,\dots,e_K)=(2,1)^{p}\,(3,1,1)^{r}
\]
has length $2p+3r=K$ and sum $3p+5r=E$. It contains neither the
factor $(2,2)$ nor the factor $(2,1,1)$, read cyclically. Let
$y_1,\dots,y_K$ be its rational cycle (Lemma~\ref{lem:ratcycle}).

\emph{Step 2: detectors and fuel on the cycle.} The parity word of
$y_i$ under $T$ (Section~\ref{sec:spine}) begins
$\mathtt O\,\mathtt E^{e_i-1}\,\mathtt O\,\mathtt E^{e_{i+1}-1}\cdots$,
and $y_i\bmod32$ is determined by its first five letters
(Terras~\cite{Terras1976}: the first $p$ letters and the residue modulo
$2^{p}$ determine each other). The residues $1$ and $9$ have parity
words $\mathtt{OEOEO}$ and $\mathtt{OEOOO}$, which would force
$(e_i,e_{i+1})=(2,2)$, resp.\ $(e_i,e_{i+1},e_{i+2})=(2,1,1)$. Neither
occurs, so no $y_i$ is $\equiv1$ or $9\pmod{32}$, and by
Lemma~\ref{lem:mod32} every $\lambda_d$ is $\bot$ on every odd integer
$\equiv y_i\pmod{32}$. Lemma~\ref{lem:fuse} holds verbatim for $2$-adic
units: a member with $e_i\ge2$ has $\tau=1$, and along a run of
$e=1$ steps $\tau$ decreases by one per step. The runs of $1$s in the
word have length at most $2$, so $\tau(y_i)\le3$ for all $i$.

\emph{Step 3: residues.} Put $M=\max(m,5)+E$. On $\Z/N'$ the maps
$\rho\mapsto(3\rho+1)2^{-e_i}$ are bijections; let $A$ be their
composite over one period and $t$ its order as a permutation. Fix
$\rho_1\in\Z/N'$ and define $\rho_{i+1}=(3\rho_i+1)2^{-e_i}$ for
$i=1,\dots,tK$ (indices of $e$ taken modulo $K$), so $\rho_{tK+1}=\rho_1$.
For $1\le i\le tK$ let $R_i$ be the residue class modulo
$R=2^{M}3^{u}N'$ of integers $x$ with
\[
x\equiv y_i\pmod{2^{M}},\qquad x\equiv y_i\pmod{3^{u}},\qquad
x\equiv\rho_i\pmod{N'}
\]
(indices of $y$ modulo $K$; the first two congruences make sense
because $y_i\in\Z_2\cap\Z_3$). Every element of $R_i$ is odd. For
$x\in R_i$: $3x+1\equiv3y_i+1\pmod{2^{M}}$, so $\vv_2(3x+1)=e_i$ and
$f(x)\equiv y_{i+1}\pmod{2^{M-e_i}}$ with $M-e_i\ge\max(m,5)$;
$3x+1\equiv3y_i+1\pmod{3^{u+1}}$, so $f(x)\equiv y_{i+1}\pmod{3^{u}}$;
and $f(x)\equiv\rho_{i+1}\pmod{N'}$. Hence for any $x\in R_i$ and
$x'\in R_{i+1}$: $f(x)\equiv x'\pmod N$; $\tau(f(x))=\tau(y_{i+1})=\tau(x')$
since $\tau(y_{i+1})\le3<M-e_i$; and all detectors are $\bot$ at both
$f(x)$ and $x'$ by Step~2. Every coordinate other than $c$ closes between
$R_i$ and $R_{i+1}$.

\emph{Step 4: cells.} Set $\delta_i=1$ if $i\bmod K\in\{1,\dots,s\}$ and
$\delta_i=0$ otherwise, so that each period carries exactly $s$ times.
Choose $c_1=H$ and $c_{i+1}=c_i+L-Qe_i+\delta_i$. Then
$c_{tK+1}=c_1+t(KL-QE+s)=c_1$. Since $e_i\le3$ and $L\ge Q$, the
downward excursion satisfies
$c_1-c_i\le t\sum_{i\le K}\max(0,Qe_i-L)\le2tKQ$.

\emph{Step 5: witnesses.} Let $\ell_0=\min(\varphi,1-\varphi)$ and take
$H\ge Q\bigl(\log_2R+j+\log_2\tfrac{2}{\ell_0\ln2}+2tK+1\bigr)$.
For each $i$, the integers in cell $c_i$ with $\theta<1-\varphi-\eta^{*}$
(if $\delta_i=0$), resp.\ $\theta\ge1-\varphi$ (if $\delta_i=1$), where
$\eta^{*}=Q/(3\cdot2^{c_i/Q}\ln2)\ge\sup_{\text{cell}}\eta$, form an
interval of length at least $2^{c_i/Q}\,\ell\ln2/Q$ with
$\ell\ge\ell_0/2$ (the choice of $H$ gives $\eta^{*}\le(1-\varphi)/2$),
which is at least $R$ by the choice of $H$ and Step~4; so the interval
contains an element $x_i$ of $R_i$. By Lemma~\ref{lem:displacement},
$c(f(x_i))=c_i+L-Qe_i+\delta_i=c_{i+1}=c(x_{i+1})$, and with Step~3 the
coordinate closes at every junction. The upward excursion is at most
$t\sum_{i\le K}\max(0,L-Qe_i+1)\le tK(L-Q+1)$, and with
$M\le m+5+E$, $E<\tfrac{8}{5}K$ and $L-Q<\tfrac35Q$ the bit length of
every witness is at most $\Theta$ when $H$ is taken minimal.

\emph{Step 6: gain.} With $y_i'=f(x_i)$,
\[
\frac{\prod y_i'}{\prod x_i}
=\frac{3^{tK}}{2^{tE}}\prod_{i}\Bigl(1+\frac1{3x_i}\Bigr)
>\Bigl(\frac{3^{K}}{2^{E}}\Bigr)^{t}>1,
\]
since $E=\lfloor K\beta\rfloor<K\beta$. Lemma~\ref{lem:certrefutes}
finishes the proof.
\end{proof}

\begin{remark}\label{rem:F}
(i) Nothing in the proof uses that the witnesses are distinct, and
no estimate of the correction terms $\eta$ beyond $\eta\le(1-\varphi)/2$
is needed: the gain is automatic from $2^{E}<3^{K}$.
(ii) Without detectors ($d=0$) any word over $\{1,2\}$ with sum $E$
serves, and then every admissible $K\ge2$ is realised, also for $j=0$;
$M$ must then be raised to $\max(m,5,2K-E+2)+E$ to cover the longer
fuel runs.
(iii) At $j=0$ and $K=2$ the word is $(2,1)$, the cycle is $-7\mapsto-5$,
and the construction reproduces the shadow of Lemma~\ref{lem:S} with its
two steps now read as two witnesses. Theorem~\ref{thm:A} is the case
$j=0$, $N=2^{m}$ of Theorem~\ref{thm:E}.
(iv) The smallest instances are tiny. At $(j,m)=(3,4)$, $K=2$:
$4123\mapsto6185$, $5817\mapsto4363$ close with
$c(6185)=c(5817)=100$, $c(4363)=c(4123)=96$, and
$6185\cdot4363>4123\cdot5817$. The pairs $x_1=2^{n}+27$, $n\ge10$,
with $x_2$ the least integer $\equiv9\pmod{16}$ in the non-carrying
part of $c(f(x_1))$ whose image returns to $c(x_1)$ with matching $\tau$,
give certificates at every height. The cost law below says why $K=2$
suffices there and why $K=7$ is needed at $j=5$.
\end{remark}

\subsection*{The cost of precision}

Theorem~\ref{thm:E} uses admissibility to construct certificates. For
witnesses that are not too small it is also necessary.

\begin{proposition}[closure criterion]\label{prop:cost}
Let $x_1,\dots,x_K$ be a closed certificate for
$\bigl(c(x),\tau(x),x\bmod2^{m}\bigr)$ with valuations
$e_i=\vv_2(3x_i+1)$ and $E=\sum e_i$, and suppose every
$x_i\ge\max\bigl(Q/(3(1-\varphi)\ln2),\ K/(3(1-\{K\beta\})\ln2)\bigr)$.
Then $QE-KL\in[0,K]$, $E\le\lfloor K\beta\rfloor$, and $K$ is admissible
at $Q$. Consequently, for $j\ge1$ and $K\ge2$, closed certificates of
length $K$ with all witnesses above any prescribed height exist if and
only if $K$ is admissible.
\end{proposition}

\begin{proof}
The first height condition gives $\eta(x_i)\le1-\varphi$, so
Lemma~\ref{lem:displacement} applies with $\delta_i\in\{0,1\}$; summing
the displacements around the cycle, $KL-QE+\sum\delta_i=0$ with
$\sum\delta_i\in[0,K]$. The gain is
$2^{-E}3^{K}\prod(1+1/(3x_i))<2^{K\beta-E}\exp\bigl(K/(3\min x_i)\bigr)$,
so gain $>1$ forces $E-K\beta<K/(3\ln2\min x_i)\le1-\{K\beta\}$ by the
second height condition, i.e.\ $E\le\lfloor K\beta\rfloor$. Then
$s(Q,K)=Q\lfloor K\beta\rfloor-KL\ge QE-KL\ge0$. The converse is
Theorem~\ref{thm:E}.
\end{proof}

\begin{proposition}[cost law]\label{prop:costlaw}
Let $K_{\min}(Q)$ be the least admissible $K$ at precision $Q$. Then
$K_{\min}(Q)\le Q$, and $\lfloor K_{\min}\beta\rfloor/K_{\min}$ is a best
approximation to $\beta$ from below among fractions with denominator at
most $K_{\min}$; hence $K_{\min}(Q)$ is the denominator of a convergent
or intermediate fraction of the continued fraction of $\log_23$ lying
below $\log_23$. For $Q=2^{j}$, $j=0,1,\dots,10$,
\[
K_{\min}=1,\,2,\,2,\,2,\,7,\,7,\,12,\,12,\,12,\,53,\,665 .
\]
\end{proposition}

\begin{proof}
$K=Q$ is admissible. If $K^{*}=K_{\min}(Q)$ then for $K<K^{*}$,
$\lfloor K\beta\rfloor/K<L/Q\le\lfloor K^{*}\beta\rfloor/K^{*}$; as
$\lfloor K\beta\rfloor/K$ is the largest fraction with denominator $K$
below $\beta$, $\lfloor K^{*}\beta\rfloor/K^{*}$ is the largest fraction
with denominator at most $K^{*}$ below $\beta$. One-sided best
approximations are exactly the convergents and intermediate fractions on
that side (Khinchin~\cite{Khinchin1964}, Ch.~I). The table is an exact
integer computation, $\lfloor n\beta\rfloor=\bitlen(3^{n})-1$.
\end{proof}

Thus the price of precision is not a continuous quantity but a Diophantine
one: it jumps at the semiconvergent denominators $2,7,12,53,665,\dots$ of
$\log_23$, and between jumps extra bits of precision are free. The
certificate miner of Section~\ref{sec:certificates}, which knows nothing
of this, finds cycles of length exactly $K_{\min}$ on the quantized-log
windows it is given ($2$ at $j=3$, $7$ at $j=5$).

\subsection*{Arbitrary moduli and multi-step potentials}

Theorem~\ref{thm:E} already allows an arbitrary modulus. The single-orbit
theorems extend the same way.

\begin{corollary}[arbitrary modulus]\label{cor:modulus}
Theorem~\ref{thm:B} remains true with $x\bmod2^{m}$ replaced by
$x\bmod N$ for any $N\ge1$. The class of Theorem~\ref{thm:A} with
$x\bmod N$ is contained in the class of Theorem~\ref{thm:E} at $j=0$.
\end{corollary}

\begin{proof}
Write $N=2^{m}3^{u}N'$ with $\gcd(N',6)=1$. On the shadow
$x=2^{N_0}w-5$ of Theorem~\ref{thm:B},
$f^{2a}(x)=\bigl(9^{a}x+5(9^{a}-8^{a})\bigr)/8^{a}$. Let $o$ be the
multiplicative order of $9\cdot8^{-1}$ modulo $N'$; if $o\mid a$ then
$9^{a}\equiv8^{a}\pmod{N'}$ and $f^{2a}(x)\equiv x\pmod{N'}$ for every
shadow $x$. If $2a\ge u$ then $9^{a}\equiv0\pmod{3^{u}}$ and
$f^{2a}(x)\equiv-5\pmod{3^{u}}$, so it suffices to take $3^{u}\mid w$.
In the proof of Theorem~\ref{thm:B} apply Dirichlet to $o\beta$ instead of
$\beta$ to obtain $a=oa'$ with $\|a\beta\|<2^{-j-3}$, and choose the
free multiplier $w$ among odd multiples of $3^{u}$ (still dense in the
logarithmic scale). The second statement is Remark~\ref{rem:F}(iii).
\end{proof}

\begin{corollary}[multi-step potentials]\label{cor:multistep}
Fix $k\ge1$. Theorems~\ref{thm:B} and~\ref{thm:E}, and hence the class of
Theorem~\ref{thm:A}, remain true for potentials that are merely
nonincreasing along $f^{k}$, i.e.\ with $\Phi(f^{k}(x))\le\Phi(x)$ for all
odd $x>1$ in place of $\Phi(f(x))\le\Phi(x)$.
\end{corollary}

\begin{proof}
For Theorem~\ref{thm:B}, apply Dirichlet to $k\beta$ to get
$a=ka'$ with $\|a\beta\|<2^{-j-3}$; the $2a$-step walk of the proof then
decomposes into $2a'$ segments of $k$ steps, and summing the $k$-step
inequalities along them gives the same contradiction. For
Theorem~\ref{thm:E}, run the proof with a word of length $K=kK_0$ and
$E=\lfloor K\beta\rfloor$, where $K_0=Q$ if $j\ge1$ and $K_0=2$ if
$j=0$ (then $2E\ge3K$ in both cases), and take as witnesses only the
segment starts $x_i\in R_{(i-1)k+1}$, $i=1,\dots,tK_0$, so that the edges
are $x_i\mapsto f^{k}(x_i)$. Over a segment with valuation sum $E_i$,
Lemma~\ref{lem:displacement} iterated gives
$c(f^{k}(x))=c(x)+\lfloor kQ\beta\rfloor-QE_i+\delta$ with
$\delta\in\{0,1\}$ decided by $\theta(x)$ against $1-\{kQ\beta\}$, and
closure needs $\sum\delta=QE-K_0\lfloor kQ\beta\rfloor$, which lies in
$[0,K_0]$ for the stated $K_0$ (it vanishes for $K_0=Q$). The rest of the
proof is unchanged, with $\{kQ\beta\}$ in place of $\varphi$ throughout.
\end{proof}

What is \emph{not} covered is noted in Section~\ref{sec:boundary}:
the leading-digit coordinate $\LD_j$ taken jointly with the quantized
logarithm, and corrections that depend on the orbit rather than on $x$.

\section{Certificates from raw orbits}\label{sec:certificates}

The constraint system ``$\Phi$ nonincreasing'', restricted to any
coordinate window, is a system of difference constraints
$g(v)-g(u)\le-\log_2(f(x)/x)$ over observed steps with
$u=\mathrm{coords}(x)$, $v=\mathrm{coords}(f(x))$; it is infeasible iff
the coordinate graph contains a directed cycle whose product of
value ratios exceeds $1$, and any such cycle is verifiable in exact
rational arithmetic. Mining by Bellman--Ford over odd $x\le10^{6}$
yields short certificates in every window we tested: the pair
$57\mapsto43$, $27\mapsto41$ closes a $2$-cycle in
$(x\bmod16,\tau)$-space with product $\tfrac{1763}{1539}>1$, refuting
Theorem~\ref{thm:A}'s class at $m\le4$ in two lines; with the level-$1$
detector added the miner returns $57\mapsto43$, $59\mapsto89$, and the
$5$-cycle $27,41,31,15,7$ through $41\mapsto31=2^{5}-1$ also closes.
On the quantized-log windows of Section~\ref{sec:quantized} with
$x\bmod16$ and $\tau$ it returns a $2$-cycle at $j=3$
($110235\mapsto165353$, $155897\mapsto116923$) and a $7$-cycle at
$j=5$, the lengths predicted by Proposition~\ref{prop:costlaw}.
Inspection of mined witnesses at larger moduli revealed the
$-5,-7$ residue pattern of Section~\ref{sec:shadow}. This pipeline is
logically independent of every lemma in this paper --- it uses raw
orbit steps only --- and would have refuted the theorems had a valid
correction existed on the tested windows.

\section{Congruence certificates on the Mersenne spine}\label{sec:spine}

The preceding sections concern potentials. This one concerns a second
method, case analysis on residue classes, applied to one explicit family
of trajectories. By Lemma~\ref{lem:fuse}, $f^{j}(2^{n}-1)=2^{n-j}3^{j}-1$
for $0\le j<n$. For odd $n\ge3$ the $n$-th step divides
$2(3^{n}-1)$ by $2^{1+\vv_2(3^n-1)}=4$ (LTE), so
\[
f^{n}(2^{n}-1)=a_n:=\frac{3^{n}-1}{2},
\]
the base-$3$ repunit, and the orbit stays above $2^{n}-1$ on the way.
The stopping time of $2^{n}-1$ (first passage below the start) is
therefore $n+\sigma(a_n)$, where $\sigma(a_n)$ is the first passage of
the orbit of $a_n$ below $2^{n}-1$. We know of no proof that
$\sigma(a_n)$ is finite for all odd $n$, let alone $O(n)$. A natural
route scores the orbit of $a_n$ block by block and case-splits on
residues of $n$. We show that such case splits are exact decision
procedures on each residue class (Proposition~\ref{prop:decide}), but
that no finite collection of them can succeed (Theorem~\ref{thm:F}).
The obstruction is again the cycle $-5\mapsto-7\mapsto-5$.

\subsection*{Blocks and scores}

Write $T(x)=x/2$ for even $x$ and $T(x)=(3x+1)/2$ for odd $x$, and read
the orbit of $x\in\Z_2$ under $T$ as a parity word in the letters
$\mathtt O$ (odd) and $\mathtt E$ (even).

\begin{definition}\label{def:block}
Let $x$ be odd with $\tau(x)=1$, so $e=\vv_2(3x+1)\ge2$, and let
$y=f(x)$. The \emph{block} at $x$ consists of the step $x\mapsto y$
followed by the $\tau(y)-1$ steps of Lemma~\ref{lem:fuse}, which end at
the next state with $\tau=1$. Its \emph{score} is
\[
\Delta=11(e-1)-9\,\tau(y).
\]
\end{definition}

In the parity word a block reads
$\mathtt O\,\mathtt E^{e-1}\,\mathtt O^{\tau(y)-1}$. It has $e-1$ even
and $\tau(y)$ odd letters, so $\Delta=11\cdot\#\mathtt E-9\cdot\#\mathtt O$.
A run of blocks with nonnegative total score has at most $11/20$ of its
letters odd. Since $3^{11}<2^{20}$, such a run contracts.
Every block begins with the factor $\mathtt{OE}$, and $\mathtt{OE}$
occurs nowhere else. Hence:

\begin{lemma}\label{lem:blockword}
The scores $\Delta_1,\dots,\Delta_B$ of the first $B$ blocks from $x$
are determined by any prefix of the parity word of $x$ that contains
the first $B+1$ occurrences of $\mathtt{OE}$.
\end{lemma}

We use the subfamily $n=64k+17$. Here $\vv_2(3^{n+1}-1)=3$ (LTE), so
$f(a_n)=x_1:=(3^{n+1}-1)/8$. Since $3^{18}\equiv9\pmod{32}$ and
$3^{64}\equiv1\pmod{256}$, $x_1\equiv1\pmod4$ and $\tau(x_1)=1$. Let
$\Delta_1,\Delta_2,\dots$ be the scores of the successive blocks from
$x_1$. For odd $k$ one has $\Delta_1=-16$ (Proposition~\ref{prop:decide}
applied to the class $k\equiv1\pmod2$). The score programme asks for
$\Delta_2+\Delta_3+\cdots\ge16$ summed over a window of $6n$ shortcut
steps. We call a class of $k$ \emph{closed at $B$} if
$\Delta_2+\cdots+\Delta_B\ge16$ on the whole class. This prefix
condition does not imply the window condition, since later blocks score
$-7,-16,\dots$ whenever $\tau(y)\ge2$. Theorem~\ref{thm:F} shows that
even the prefix condition cannot be certified on all odd $k$.

\subsection*{Exact decision on a residue class}

\begin{lemma}\label{lem:determinacy}
For $P\ge6$, $x_1\bmod2^{P}$ depends only on $k\bmod2^{P-5}$.
Moreover, if $x\equiv x'\pmod{2^{p}}$ in $\Z_2$, then the parity words
of $x$ and $x'$ agree in their first $p$ letters.
\end{lemma}

\begin{proof}
The order of $3$ modulo $2^{P+3}$ is $2^{P+1}$, so
$3^{64k+18}\bmod2^{P+3}$ depends on $64k\bmod2^{P+1}$. The second
statement is Terras's observation~\cite{Terras1976}: $T^{i}(x)\equiv
T^{i}(x')\pmod{2^{p-i}}$ for $i\le p$, by induction.
\end{proof}

\begin{proposition}[exact class decision]\label{prop:decide}
Fix a class $k\equiv r\pmod{2^{j}}$. The first $j+5$ letters of the
parity word of $x_1$ are the same for every $k$ in the class, including
every $2$-adic $k$. If they contain $B+1$ occurrences of $\mathtt{OE}$,
then $(\Delta_1,\dots,\Delta_B)$ is constant on the class. Otherwise
split the class into its two children modulo $2^{j+1}$ and repeat. A
statement about $(\Delta_1,\dots,\Delta_B)$ on a class is proved when it
holds on every leaf of this tree, and refuted, with an explicit integer
counterexample, when it fails on some leaf.
\end{proposition}

\begin{proof}
Lemmas~\ref{lem:determinacy} and~\ref{lem:blockword}.
\end{proof}

For example, single class evaluations give
\[
\begin{array}{c|c|c}
k\ \text{class} & (\Delta_2,\Delta_3,\Delta_4) & \text{sum}\\\hline
3\bmod2^{8} & (2,13,2) & 17\\
171\bmod2^{8} & (2,2,13) & 17\\
323\bmod2^{9} & (2,13,13) & 28\\
579\bmod2^{10} & (2,13,4) & 19
\end{array}
\]
so these four classes are closed at $B=4$.

\begin{remark}\label{rem:onerep}
A single representative is not a certificate. The integers $k=267$ and
$k=41227$ lie in the same class modulo $2^{13}$. For $k=267$,
$\Delta_2+\cdots+\Delta_6=43$. For $k=41227$ the same sum is $7$. The
class modulo $2^{13}$ does not determine $\Delta_6$.
Proposition~\ref{prop:decide} refines it, and the leaf containing $41227$
has modulus $2^{18}$.
\end{remark}

\subsection*{The cycle defeats every finite certificate}

\begin{lemma}\label{lem:ball}
The map $k\mapsto x_1(k)$ extends to a homeomorphism from the odd
$2$-adic integers onto the ball $9+64\Z_2$.
\end{lemma}

\begin{proof}
Write $3^{64}=1+2^{8}u$ with $u$ odd. Then $z\mapsto(1+2^{8}u)^{z}$
is a homeomorphism from $\Z_2$ onto $1+2^{8}\Z_2$ that scales distances
by exactly $2^{-8}$, and odd $z$ go onto $257+2^{9}\Z_2$. The affine map
$w\mapsto(3^{18}w-1)/8$ takes $257+2^{9}\Z_2$ onto
$(3^{18}\cdot257-1)/8+64\Z_2=9+64\Z_2$, using
$3^{18}\cdot257\equiv73\pmod{512}$.
\end{proof}

So every point of $9+64\Z_2$ is $x_1$ of some odd $2$-adic $k$. One of
them falls straight into the expanding cycle of
Section~\ref{sec:shadow}.

\begin{theorem}[shadow of the $-5$ cycle on the spine]\label{thm:F}
Let $k^{*}\in\Z_2$ be defined by $3^{64k^{*}+20}=-367$. Then
\begin{enumerate}
\item $k^{*}$ is odd, $x_1(k^{*})=-47/9$, and its blocks have
  $\Delta_1=-16$ and $\Delta_j=-7$ for every $j\ge2$;
\item for every $J\ge1$ and every integer $k\ge0$ with
  $k\equiv k^{*}\pmod{2^{J}}$,
  \[
  \Delta_2=\Delta_3=\cdots=\Delta_{B}=-7,
  \qquad B=\Bigl\lfloor\frac{J+2}{3}\Bigr\rfloor;
  \]
\item no finite set of residue classes of $k$, each closed at some $B$
  by Proposition~\ref{prop:decide}, contains every odd integer $k$.
\end{enumerate}
\end{theorem}

\begin{proof}
(1) One has $-367\cdot3^{-20}\equiv257\pmod{512}$, so $k^{*}$ exists
and is odd by the proof of Lemma~\ref{lem:ball}. Then
$x_1=(3^{64k^{*}+18}-1)/8=(-367/9-1)/8=-47/9$. Next
$3x_1+1=-44/3$, so $e=2$ and $y_1=-11/3$ with
$\tau(y_1)=\vv_2(-8/3)=3$, giving $\Delta_1=11-27=-16$. Two steps of
Lemma~\ref{lem:fuse} give $-11/3\mapsto-5\mapsto-7$. From $-7$ every
block is $-7\mapsto-5\mapsto-7$ with $e=2$ and $\tau(-5)=2$, so
$\Delta=11-18=-7$.

(2) The parity word of $-47/9$ is
$\mathtt{OEOO}(\mathtt{OEO})^{\infty}$. Its $(B+1)$-st occurrence of
$\mathtt{OE}$ ends at letter $3B+3$. If $k\equiv k^{*}\pmod{2^{J}}$,
Lemma~\ref{lem:determinacy} gives agreement in the first $J+5$ letters,
which contain $B+1$ occurrences of $\mathtt{OE}$ when $3B+3\le J+5$.
Lemma~\ref{lem:blockword} finishes.

(3) Suppose classes $C_1,\dots,C_m$ cover the odd integers. Refine them
to a common modulus $2^{J}$. Every odd residue modulo $2^{J}$ is then
covered, including that of $k^{*}$, so some $C_i$ contains $k^{*}$. The
certificate for $C_i$ uses only letters common to every $2$-adic member
of $C_i$ (Proposition~\ref{prop:decide}), so it applies to $k^{*}$ and
gives $\Delta_2+\cdots+\Delta_{B_i}\ge16$. By (1) that sum is
$-7(B_i-1)<16$, a contradiction.
\end{proof}

\begin{remark}[relation to Terras]\label{rem:terras}
Terras~\cite{Terras1976} showed that first passage below the starting
point within $p$ steps is a property of $x\bmod2^{p}$, and the burn of
Lemma~\ref{lem:fuse} shows that the class $x\equiv-1\pmod{2^{p}}$ has
not stopped by step $p$, so no finite modulus decides the stopping time
of all integers. Theorem~\ref{thm:F} is the corresponding statement for
the \emph{index} $n$ of the thin family $2^{n}-1$, where residue classes
of $n$ are a different and a priori more promising partition; it
identifies the obstruction as one explicit $2$-adic index and one
explicit rational cycle, and it concerns block scores, which are finer
than stopping.
\end{remark}

The mechanism is the one behind Theorem~\ref{thm:A}. The cycle block
$-7\mapsto-5\mapsto-7$ has one even and two odd letters and value ratio
$3^{2}/2^{3}=\tfrac98>1$. Any linear score whose nonnegativity forces
contraction must therefore score this block negatively. Residue classes
of $k$ cannot separate $k$ from $k^{*}$, just as residues of $x$ cannot
separate $x$ from its shadow. In (2), each extra $3$ bits of modulus
buy exactly one more block of shadowing. One checks that $B$ is sharp
for $J\le120$.

\subsection*{How far certificates reach}

Proposition~\ref{prop:decide} also measures the method. Among odd $k$,
the Haar measure of the classes modulo $2^{J}$ that are closed within
ten blocks is
\[
\frac{1013}{2048}\approx0.495\ (J=14),\qquad
\frac{77177}{131072}\approx0.589\ (J=18),\qquad
\frac{10687}{16384}\approx0.652\ (J=22).
\]
The rest is undetermined at that modulus, apart from a part of measure
about $2.5\times10^{-4}$ at $J=22$ that provably does not close within
ten blocks. Each four extra bits close about $6\%$ more. By
Theorem~\ref{thm:F}(3) the measure never reaches $1$ at any finite $J$.

\begin{remark}\label{rem:spinebarrier}
Even a complete cover would yield only prefix inequalities. The window
condition also needs the later blocks, and a lower bound for them is the
same problem for the shifted orbit. The full score programme of the
repository accompanying this paper, if completed, would give
$\sigma(a_n)\le6n$ and hence a stopping time of at most $7n$ odd steps
for every $2^{n}-1$ with $n$ odd. That is a linear stopping-time bound for an
explicit infinite family whose orbits first climb by a factor of about
$(3/2)^{n}$, and we know of no such result.
\end{remark}

\section{The boundary}\label{sec:boundary}

Theorems~\ref{thm:A}, \ref{thm:B} and~\ref{thm:E}, with
Corollaries~\ref{cor:modulus} and~\ref{cor:multistep}, exclude every
correction of the form
\begin{align*}
&g\bigl(x\bmod N,\ \tau(x),\ \lambda_1(x),\dots,\lambda_d(x),\
\lfloor2^{j}\log_2x\rfloor\bigr)\\
\text{or}\quad
&g\bigl(x\bmod N,\ \tau(x),\ \lambda_1(x),\dots,\lambda_d(x),\
\LD_j(x)\bigr),
\end{align*}
for all $N,d,j$, along $f^{k}$ for every fixed $k$. Every exclusion is
witnessed by a finite closed cycle in a coordinate graph, every cycle is
a shadow of an expanding rational cycle, and the one quantitative
invariant that appears --- the number of witnesses needed at precision
$2^{j}$ --- is the Diophantine cost of approximating $\log_23$ from
below. Two things are not covered.

\begin{question}\label{q:joint}
Is there a nonincreasing potential
$\log_2x+g\bigl(\lfloor2^{j}\log_2x\rfloor,\ \LD_{j'}(x),\ \tau(x),\
x\bmod2^{m}\bigr)$ for some $j,j',m$? The two size coordinates have
breakpoints at $2^{k/2^{j}}$ and at $\ell\,2^{\bitlen(x)-j'}$, which do
not nest, and the construction of Theorem~\ref{thm:E} fixes the position
of each witness in its cell only up to the width of a sub-interval. We
expect a rotation argument to handle the joint coordinate, but have not
carried it out. More generally one may ask the same for any finite
partition of the mantissa $x/2^{\bitlen(x)-1}\in[1,2)$.
\end{question}

\begin{question}\label{q:path}
Do path-dependent accountings --- functions of the orbit segment leading
to $x$ rather than of $x$ --- admit analogous obstructions? Every
coordinate treated here is a function of $x$ alone, and the certificates
exploit exactly that: a witness may be chosen afresh at every junction.
A correction that reads how $x$ was reached is a different object, and
as the memory grows the existence of a monotone potential becomes
conjecture-adjacent.
\end{question}

We also note what the results do \emph{not} say. Apart from the
negative statements of Section~\ref{sec:spine} about one method on one
family, nothing here bears on individual trajectories, densities of
divergent orbits, or cycles. The theorems delimit a proof \emph{method},
not the dynamics.

\section{Computational verification}\label{sec:verification}

{\sloppy
Every numbered claim is verified by scripts using exact integer and
rational arithmetic (no floating point in any assertion): the shadow
lemma over wide parameter ranges with detector evaluation at all nodes
(\texttt{verify\_shadow.py}); the iterated shadow, the concrete
Dirichlet instance $(m,j,a)=(6,4,53)$ at depth $N=177$, and the
irrationality arithmetic (\texttt{verify\_leading\_digit.py}); the
tower in all parts, including both directions of the domain criterion
and the carry-free normal form
(\texttt{verify\_tower.py}); the family-method subclasses
(\texttt{verify\_no\_local\_potential.py}, \texttt{verify\_nlp2.py});
the independent certificate miner with exact cycle confirmation
(\texttt{verify\_nogo\_certificate.py},
\texttt{verify\_fuel\_fraction.py}); for Section~\ref{sec:quantized},
the rational-cycle lemma over all $21\,844$ words with $K\le7$,
$e_i\le4$, the construction of Theorem~\ref{thm:E} on a grid of
$(j,K,N,d)$ with $j\le8$ (closure at every junction, detectors,
arbitrary moduli including $N=2^{5}\cdot3^{2}\cdot5\cdot7$, and the
exact product gain), the admissibility table, the modulus and multi-step
corollaries, and the explicit instances of Remark~\ref{rem:F}
(\texttt{verify\_quantized\_log\_theorem.py}); and for Section~\ref{sec:spine}, the
ball lemma, the shadow index $k^{*}$, part~(2) of Theorem~\ref{thm:F} for
$J\le120$, the four class evaluations, and the exact closure measures
(\texttt{verify\_residue\_certificate\_nogo.py --measure-22}), with the
class decision procedure in \texttt{prove\_block8\_17\_classes.py}. The
scripts accompany the repository.
\par}

\begin{thebibliography}{99}

\bibitem{Andaloro2000}
P.~Andaloro,
\emph{On total stopping times under $3x+1$ iteration},
Fibonacci Quart. \textbf{38} (2000), 73--78.

\bibitem{Barina2021}
D.~Barina,
\emph{Convergence verification of the Collatz problem},
J. Supercomputing \textbf{77} (2021), 2681--2688.

\bibitem{Everett1977}
C.~J. Everett,
\emph{Iteration of the number-theoretic function $f(2n)=n$,
$f(2n+1)=3n+2$},
Adv. Math. \textbf{25} (1977), 42--45.

\bibitem{Khinchin1964}
A.~Ya. Khinchin,
\emph{Continued Fractions},
University of Chicago Press, Chicago, 1964.

\bibitem{KrasikovLagarias2003}
I.~Krasikov and J.~C. Lagarias,
\emph{Bounds for the $3x+1$ problem using difference inequalities},
Acta Arith. \textbf{109} (2003), 237--258.

\bibitem{Lagarias1985}
J.~C. Lagarias,
\emph{The $3x+1$ problem and its generalizations},
Amer. Math. Monthly \textbf{92} (1985), 3--23.

\bibitem{Lagarias2010}
J.~C. Lagarias (ed.),
\emph{The Ultimate Challenge: The $3x+1$ Problem},
Amer. Math. Soc., Providence, RI, 2010.

\bibitem{LagariasBibI}
J.~C. Lagarias,
\emph{The $3x+1$ problem: an annotated bibliography (1963--1999)},
arXiv:math/0309224.

\bibitem{LagariasBibII}
J.~C. Lagarias,
\emph{The $3x+1$ problem: an annotated bibliography, II (2000--2009)},
arXiv:math/0608208.

\bibitem{SimonsDeWeger2005}
J.~Simons and B.~de Weger,
\emph{Theoretical and computational bounds for $m$-cycles of the $3n+1$
problem},
Acta Arith. \textbf{117} (2005), 51--70.

\bibitem{Steiner1977}
R.~P. Steiner,
\emph{A theorem on the Syracuse problem},
Proc. 7th Manitoba Conference on Numerical Mathematics, 1977,
553--559.

\bibitem{Tao2022}
T.~Tao,
\emph{Almost all orbits of the Collatz map attain almost bounded
values},
Forum Math. Pi \textbf{10} (2022), e12.

\bibitem{Terras1976}
R.~Terras,
\emph{A stopping time problem on the positive integers},
Acta Arith. \textbf{30} (1976), 241--252.

\bibitem{Wirsching1998}
G.~J. Wirsching,
\emph{The Dynamical System Generated by the $3n+1$ Function},
Lecture Notes in Math. 1681, Springer, Berlin, 1998.

\end{thebibliography}

\end{document}
